\section*{Note: Emergent Fermionic Statistics and Accessibility Barriers in the Hilbert Substrate}

\subsection*{Summary of Findings}

In this work we investigated the emergence of matter-like and spacetime-like structure in the \emph{Hilbert Substrate Framework} (HSF) using small-$N$ numerical experiments on XX-type spin-chain substrates. The central question was whether fermionic statistics and spatial locality arise as independent or sequential features of the substrate dynamics.

Our simulations establish three empirical findings:

\begin{enumerate}
    \item \textbf{Emergent fermionic statistics.}  
    Across multiple seeds and scrambling regimes, the accessible operator algebra admits a fermionic effective description. Jordan--Wigner--type mode operators satisfy anticommutation relations to numerical tolerance, sector energies exhibit additive behavior consistent with noninteracting fermions, and exclusion-like ``Pauli pressure'' diagnostics appear as curvature in the accessible landscape. These signatures arise without assuming geometry \emph{a priori}.

    \item \textbf{Accessibility-limited locality.}  
    Spatial locality is not a basis-invariant property of the substrate. When the Hamiltonian is scrambled by products of local (single-site or low-depth) unitaries, locality-biased recovery procedures reliably restore nearest-neighbor structure. In contrast, when the Hamiltonian is scrambled by generic global unitaries, locality is not recoverable using local generators alone. This behavior is robust across diagnostics and implementations.

    \item \textbf{Persistence of matter without geometry.}  
    Crucially, fermionic diagnostics remain intact even in regimes where spatial locality is dynamically inaccessible. Global scrambling destroys the path back to a spatially local basis, yet exchange statistics and sector structure persist. This demonstrates that matter-like structure can precede, and does not require, emergent spacetime locality.
\end{enumerate}

\subsection*{Interpretation}

These results support an ordering of emergence in which
\[
\text{Hilbert substrate}
\;\rightarrow\;
\text{subsystems}
\;\rightarrow\;
\text{exchange statistics (matter)}
\;\rightarrow\;
\text{accessibility constraints}
\;\rightarrow\;
\text{locality}
\;\rightarrow\;
\text{geometry}.
\]

In this picture, spatial locality occupies a restricted accessibility basin in Hilbert space rather than representing a generic or recoverable feature under arbitrary basis changes. Global scrambling constitutes a transition across an accessibility barrier, beyond which locality cannot be restored using local operations alone. The irreversibility of this transition is itself the signal.

\subsection*{Scope and Limitations}

The present results are obtained in small-$N$ toy systems and do not constitute a derivation of continuum spacetime or macroscopic dimensionality. Rather, they establish that fermionic exchange statistics emerge robustly within the substrate dynamics, while locality and geometry are contingent, fragile, and accessibility-limited. Larger systems and extended substrates will be required to investigate dimensional selection and continuum limits.

\subsection*{Conclusion}

Within the tested regime, fermionic statistics emerge as a stable and robust feature of the Hilbert Substrate, whereas spatial locality appears as a secondary, dynamically constrained phase. The failure to recover locality after global scrambling is not a defect of the framework but a confirmation of its central claim: spacetime structure is emergent, contingent, and restricted by accessibility rather than assumed from the outset.
